\section{Data}

Our analysis relies on data scraped from Twitter (X). We measure behavioral outcomes by first identifying a sample of followers with public profiles for each influencer. This sample includes followers who have very strong or strong ties with the influencer (as defined above), as well as a random sample of followers with weak ties to the influencer. Due to restrictions imposed by the Twitter API, we are limited in our ability to track all weak ties. We monitor these followers' posts and social engagement (%likes, 
shares and tweets) during and after the study. 
We use also measure whether followers follow Africa Check and social media influencers.

To enhance the precision of our estimates, we also filter out followers at both extremes of the engagement distribution. Baseline information was employed to identify and exclude potential bots, which we define as users posting an anomalously high volume of tweets within a short period \footnote{Users above the 95th percentile of tweets posted at baseline were excluded.}. Moreover, we also filter out followers who did not post any tweets during their baseline periods, thereby ensuring our analysis concentrated on active users. After aggregating followers from both South Africa and Kenya, the final sample for the intensive margin analysis consists of 169,951 followers.

Our analysis is organized into four subcategories: verifiability, COVID sentiment, URLs, and interactions. For each subcategory, we define and, in some cases, create various outcome variables that enable us to assess the impact of the intervention on users' overall behavior.

\textbf{Activity analysis}

\textbf{Verifiable analysis} %To understand the effects of our intervention on the behavior of the followers when interacting or sharing misleading content, we analyzed all English tweets posted or retweeted by them. Given the vast number of posts by followers, 
We trained a Machine Learning model to label posts as verifiable or. Specifically, thanks to Africa Check' help in labeling a sample of posts, we trained a binary classification model based on the BERT architecture. This model labeled posts as either verifiable or non-verifiable, and had a validation accuracy of 85\%.\footnote{Refer to Appendix \ref{ver_bert} for a detailed description of the model, the training dataset and the validation accuracy.} Beyond being of interest of itself, by identifying verifiable posts, we could subsequently label verifiable posts as fake or true. 

\textbf{Fake/true analysis}
%Importantly, we were particularly interested in the change in behavior related to misleading claims or actual misinformation. To analyze this, 
We developed an additional model to label verifiable posts as either fake or true. This binary classification model was trained using posts labeled as ``fake'' by Africa Check in their fact checks or ``true'' from reputable news sources, such as established newspapers. Similar to our verifiable model, we utilized the BERT architecture. Our model had a validation accuracy of 84\%, \footnote{Refer to Appendix \ref{fake_bert} for a detailed description of the model, the training dataset and the validation accuracy. And to Appendix \ref{lda_fake} for a comprehensive description of how we built the training data.} 
%The primary objective of this model was to predict whether scraped posts and tweets contained false information\footnote{Or if they were at least more similar to a misinformation post than to a truthful one.}

\textbf{COVID-19 sentiment analysis} %Given that some of the information shared by social media micro-influencers and Africa Check through Twitter (X) Ads relates to COVID-19 and vaccination campaigns, we decided to 
We analyze the effect of our intervention on the sentiment expressed by followers towards COVID-19 and the COVID-19 vaccines. To accomplish this, we first identified and labeled tweets containing any information related to COVID-19 and COVID-19 vaccines. We manually defined a set of keywords, such as ``coronavirus'', ``COVID-19'', ``mRNA'', ``Pfizer'', etc., and filtered all posts containing these terms. From the subset of posts containing these keywords, we extracted the most frequent words (excluding the initial set of keywords) to ensure that no relevant terms related to the topic were excluded. We then identify the set of posts containing any of these terms.

We employ two different sentiment analysis models to predict the sentiment of each post. Our primary analysis utilizes a BERT model, trained on a comprehensive dataset sourced from Kaggle, to perform sentiment prediction on the tweets.\footnote{Appendix \ref{app:covid_bert} provides a detailed description of the dataset, the model used, and the evaluation metrics.} Additionally, for robustness, we consider \href{https://vadersentiment.readthedocs.io/en/latest/}{VADER}, which is is a lexicon and rule-based sentiment analysis tool specifically designed to detect sentiments expressed on social media. Specifically, we label posts as positive, negative, or neutral.
%we aimed to determine the number of positive, negative, and neutral tweets related to these topics posted by each influencer at each stage of the study. Therefore, for each follower we generate the total number, as well as the positive, neutral and negative Tweets (posts and RTs) that refer to COVID or also, COVID vaccines.

\textbf{URL analysis} %Additionally, knowing that most 
Since information and misinformation are often not explicitly written in posts, but shared through links, we analyze the URLs shared by followers. To do this, we first extracted all URLs from the posts, grouped them by domain name,\footnote{The net domain name for famous sites such as Daily Mail UK and Fox News are ``www.dailymail.co.uk'' ``www.foxnews.com'' respectively} and manually classified them into the following categories: news articles from from reliable and non-reliable sources, fact checks,  other information sources such as blogs, and other websites that do not provide information, such as gambling websites.

%We generated seven different outcomes from this analysis: the total number of URLs shared by each user in a given period, the total number of URLs linking to news articles, the total number of news articles from reliable and non-reliable sources, the total number of fact-checks shared, other information sources such as blogs and NGOs, and other websites such as gambling sites (classified as Not Info.).

To categorize the news websites from Africa, we first leveraged a dataset from Africa Check classifying reliable and non-reliable new websites. Additionally, Africa Check reviewed and classify an extra batch of African news websites. For global news websites, we used several sources, including \href{https://www.newsguardtech.com/}{NewsGuard} and \href{https://mediabiasfactcheck.com/pro-science/}{MediaBias}.\footnote{Her we share 4 examples of sources we used to categorize the newsites, most coming from NewsGuard. For reliable newsites we encountered several lists from reliable news sources such as \href{https://www.forbes.com/sites/berlinschoolofcreativeleadership/2017/02/01/10-journalism-brands-where-you-will-find-real-facts-rather-than-alternative-facts/?sh=2ae983c2e9b5}{Forbes} and also directly from \href{https://web.archive.org/web/20230110151501/https://www.newsguardtech.com/special-reports/uk-special-report-the-best-and-worst-of-2022/}{NewsGuard}. For non-reliable we used mostly NewsGuard's veracity grading articles, for example for \href{https://www.newsguardtech.com/press/newsguard-launches-russia-ukraine-disinformation-tracking-center/}{RT} and for \href{https://www.newsbusters.org/blogs/free-speech/catherine-salgado/2022/04/28/whoa-newsguard-finally-downgrades-daily-beast-hunter}{The Daily Beast Hunter}.}

\textbf{Interactions} Before the change in pricing of the Twitter API, we could freely access interactions between each post, including the number of likes, retweets, comments, and quotes each tweet received. We use this data to understand whether the popularity of each post that followers interacted with.\footnote{We focus exclusively on original tweets, not retweets, as the API provides interaction metrics for the initial tweet in the case of retweets.} %Specifically, we aimed to observe changes in interactions for labeled tweets, such as the average change in likes that treated users received for their verifiable posts compared to control users.

To conduct this analysis, we defined consider the total interactions (likes, comments, and retweets) with their posts, as well as those specific to different types of posts such as verifiable and non-verifiable posts, and posts with different sentiments about COVID-19 and the COVID-19 vaccines.
